\section{Discussion}
\label{sec:discussion}

The central aim of this thesis was to investigate whether a lightweight neuro-symbolic anomaly-analysis framework could provide a more interpretable, operationally meaningful, and deployment-conscious alternative to purely statistical anomaly detection for weather-dependent electricity systems. More specifically, the research sought to determine whether a compact dCeNN--ELM forecasting backbone, when combined with calibrated residual analysis and ASP-based symbolic refinement, could improve the plausibility and usefulness of anomaly signals under multivariate uncertainty. The results developed in the preceding chapters suggest that the answer is \emph{qualified but positive}: the proposed framework does not emerge as the strongest pure forecaster in all respects, but it does demonstrate value in the dimensions that were most central to the thesis, namely interpretability, modularity, symbolic plausibility control, and edge-oriented deployability \cite{abeysekara2025report,repo_benchmark_comparison,repo_metrics_monthly,repo_veto_analysis_plot,repo_edge_export_ch62}.

This nuanced outcome is important. If the thesis were framed as a claim of universal predictive superiority, the reported results would not fully support that position, particularly given the strong performance of the Ridge baseline in the committed benchmark outputs \cite{repo_benchmark_comparison}. However, the actual contribution of the research is broader and more system-oriented than that. The framework is designed as a full anomaly-analysis pipeline rather than as a single forecasting contest entry. It includes structured data engineering, lightweight latent forecasting, residual-based candidate generation, symbolic veto reasoning, finance-aware interpretation, and edge export. The discussion therefore evaluates the work not only by raw forecast fit, but by how well it answers the research questions and fulfills the original objectives of producing a lightweight, interpretable, and practically useful neuro-symbolic system.

At a high level, the thesis objectives can be interpreted as three linked goals:
\begin{enumerate}
	\item to develop a forecasting backbone that is sufficiently lightweight for CPU-first and edge-oriented deployment;
	\item to improve anomaly interpretability and false-alarm control through symbolic reasoning;
	\item to move anomaly evaluation beyond raw counts toward operational or financial usefulness.
\end{enumerate}
The discussion that follows argues that the first objective is \emph{partially achieved}, the second is \emph{substantially achieved}, and the third is \emph{promisingly addressed but best interpreted with appropriate caution}. This pattern is scientifically acceptable and, in fact, characteristic of good thesis work: it shows where the framework is strongest, where it remains mixed, and where its results motivate future development.

\subsection{Interpretability}
\label{subsec:discussion_interpretability}

Interpretability is one of the clearest strengths of the proposed framework. In conventional residual-based anomaly detection, a flagged timestamp is usually justified only by the magnitude of a residual or anomaly score. While this can indicate that the observation is statistically unusual, it does not explain \emph{why} the signal should be accepted as meaningful or \emph{why} it should be rejected as implausible. The ASP component addresses this limitation directly by introducing explicit logical criteria for anomaly acceptance and rejection \cite{gelfond1988stable,gebser2012answer,gebser2019multi}.

The key interpretability gain comes from the fact that the symbolic layer transforms anomaly evaluation from a purely numerical decision into a rule-grounded one. A raw anomaly can be vetoed because it violates a physical plausibility constraint, because it lacks persistence, because it coincides with calendar context that weakens its anomaly interpretation, or because it lacks the cross-signal confirmation required by the market logic. This is fundamentally different from post-hoc feature attribution in black-box models. The explanation is not merely that a variable was influential; rather, the explanation is that the event failed a human-readable domain constraint. In methodological terms, this gives the anomaly detector a \emph{semantic rejection criterion}, not only a statistical threshold.

The results chapter already indicated how this interpretability becomes operational. The veto reasoning analysis and the report-level veto summaries provide explicit categories for rejected anomalies, such as night-time solar inconsistency, wind-ramp plausibility violations, holiday or weekend context exemptions, missing co-anomaly support, and transient non-persistent spikes \cite{repo_veto_analysis_plot,repo_report_tables}. Even where the final report tables use proxy-style reasoning categories rather than raw solver-emitted natural-language text, the interpretive value remains substantial. Each rejected anomaly can be associated with a recognizable domain logic, which is far more informative than simply labeling it as a ``false positive.''

This directly addresses one of the core research questions of the thesis: whether anomaly detection can be made more interpretable without abandoning data-driven learning. The answer supported by the present work is yes. The ASP layer provides a principled ``why'' for anomaly rejection, and therefore closes an important gap between black-box statistical detection and domain-grounded reasoning. In this sense, the research objective of producing an interpretable anomaly-analysis pipeline is convincingly satisfied.

At the same time, interpretability in the current framework should not be overstated. The rule base is intentionally compact, and therefore the symbolic explanations are limited to the subset of domain logic that has been encoded. This means that the system explains anomalies in terms of the \emph{implemented} rules, not in terms of a complete model of electricity-market behaviour. Nevertheless, this is still a meaningful advance over purely opaque anomaly scoring. The thesis does not claim exhaustive explainability; it demonstrates a credible and extensible path toward rule-grounded anomaly interpretation.

\subsection{Scalability}
\label{subsec:discussion_scalability}

Scalability presents a more nuanced picture. On the one hand, the software architecture and the forecasting backbone are both deliberately designed for lightweight operation. The dCeNN--ELM model is computationally modest, the parameter count is small relative to the LSTM baseline, the inference latency is low, and the deployment path is simplified through a NumPy-only runtime export \cite{repo_benchmark_comparison,repo_edge_export_ch62,repo_readme_ch62}. These are strong indicators that the neural component of the framework scales well in the sense most relevant to the thesis, namely CPU-first execution on modest hardware.

On the other hand, the symbolic layer introduces a different kind of scaling challenge. ASP reasoning is highly attractive for interpretability and rule consistency, but its cost depends not only on the number of anomalies, but also on the number of grounded facts, the temporal window over which rules are applied, and the combinatorial complexity of the rule base \cite{gebser2012answer,gebser2019multi}. In the present thesis, this challenge is moderated by several factors: the time resolution is hourly rather than sub-second, the rule set is compact, the number of signals is limited, and symbolic reasoning is applied as a post-detection refinement layer rather than as a continuously firing online theorem prover. Under these conditions, the approach is quite reasonable.

However, if one imagines scaling the same architecture to very high-frequency streams, very large multi-zone systems, or much richer rule bases, the trade-off becomes sharper. The neural component would still remain relatively manageable, but the symbolic refinement stage could become the bottleneck. This does not invalidate the approach; rather, it clarifies its current operating regime. The present framework is best understood as suitable for \emph{hourly or near-real-time} monitoring in moderate-scale environments, especially where interpretability matters enough to justify some additional reasoning overhead.

This leads to an important conclusion regarding the corresponding research objective. If scalability is understood as ``can the system operate as an edge-conscious hourly anomaly-analysis pipeline on modest CPU resources?'', then the research gives a positive answer. If scalability is understood more ambitiously as ``can the same ASP-enhanced reasoning stack scale without modification to large, high-frequency, real-time infrastructures?'', then the answer is more cautious. The framework is scalable enough for the thesis setting and its intended deployment narrative, but full real-time large-scale symbolic reasoning would likely require additional engineering strategies such as sliding-window grounding, incremental solving, rule pruning, two-stage screening, or approximate reasoning layers.

This trade-off is intellectually valuable rather than problematic. It shows that the thesis does not merely assert neuro-symbolic appeal in the abstract; it also reveals the engineering boundary at which symbolic expressiveness begins to compete with real-time constraints. That is precisely the kind of honest systems discussion expected in strong research work.

\subsection{Robustness under Weather Uncertainty}
\label{subsec:discussion_weather_uncertainty}

Robustness under weather uncertainty is one of the most important thematic tests of the framework, because the chosen application domain is explicitly weather-dependent. The Austrian 2022 data include periods of pronounced seasonal and market instability, and the forecasting analysis already showed that model performance is not temporally uniform \cite{repo_metrics_monthly}. In particular, summer months---especially August---were among the most difficult for several targets, especially price and wind, with visibly elevated forecast errors and stronger anomaly burden. These findings indicate that the system is indeed being tested under conditions where uncertainty is real and not artificially simplified.

From a strict predictive standpoint, the results show that the framework is not perfectly robust in the sense of maintaining uniformly strong forecast fit under all weather regimes. That conclusion must be stated clearly. The seasonal degradation in price and wind forecasting, together with the negative $R^2$ values reported in the detailed forecasting summaries, indicates that the neural backbone remains vulnerable to difficult and volatile periods \cite{repo_metrics_summary,repo_metrics_monthly}. This means that weather uncertainty is not ``solved'' by the forecasting model alone.

Yet this is precisely where the layered design of the thesis becomes important. Robustness in the present work should not be interpreted only as stable RMSE. It should be interpreted as the \emph{ability of the overall system to remain analytically useful when weather-driven volatility increases}. On that criterion, the framework performs better. Feature engineering injects weather-sensitive capacity factors, radiation, wind proxies, and calendar context into the model input. Thresholding is calibrated in a context-aware manner. The symbolic layer then filters residual-based alerts through plausibility constraints that are especially relevant in weather-sensitive conditions, such as solar inconsistency under zero radiation and persistence checks for wind-related anomalies. Finally, the case-study framework shows that the system can still isolate meaningful high-severity events during difficult periods, rather than collapsing into uninterpretable alert noise \cite{repo_anomaly_events_2022,repo_plot_master_timeline,repo_additional_visuals}.

This layered response to uncertainty is a major strength of the thesis. Instead of assuming that one model class alone will be robust enough to handle all meteorological variability, the framework distributes the burden: forecasting handles expected structure, thresholding handles uncertainty, ASP handles plausibility, and case-based analysis plus utility evaluation handle practical interpretation. That is a more realistic approach to weather-dependent anomaly analysis than simply trying to force all robustness into the neural model itself.

The selected high-severity case studies reinforce this interpretation. The strongest price and wind events in the committed event table occur precisely in late-summer periods where forecasting is hardest. The fact that these events remain visible, structured, and analyzable within the event-level pipeline suggests that the framework retains practical usefulness even when prediction error rises. In this sense, the research objective of handling anomaly analysis under weather uncertainty is met in a \emph{qualified but meaningful} way: the system does not eliminate weather-induced forecasting difficulty, but it does remain capable of surfacing and refining operationally important events in the presence of that difficulty.

\subsection{Overall Assessment Against the Research Questions and Objectives}

Taken together, the results support a balanced answer to the overarching research questions. The thesis asked, in effect, whether a lightweight neuro-symbolic pipeline could improve anomaly analysis in three ways: by making it more interpretable, by keeping it computationally practical, and by making it more useful under real-world uncertainty. The discussion above suggests the following answer.

First, the framework \textbf{does succeed in improving interpretability}. This is one of its strongest contributions. The ASP-based veto layer gives rule-grounded reasons for rejecting anomaly candidates, thereby providing a level of semantic transparency that purely residual-based detectors lack.

Second, the framework \textbf{is practically lightweight at the neural level and reasonably scalable at the full pipeline level}, but with clear trade-offs. The dCeNN--ELM backbone supports the thesis claim of compact forecasting and edge-feasible inference, while the ASP stage remains appropriate for hourly and moderate-scale reasoning rather than unrestricted real-time expansion.

Third, the framework \textbf{handles weather uncertainty in a layered rather than monolithic way}. It does not eliminate all forecasting difficulty under volatile meteorological or market conditions, but it remains analytically useful by combining context-aware features, calibrated residual logic, symbolic refinement, and event-level interpretation.

This means that the thesis objectives are not answered in a simplistic ``all metrics improved'' sense. Instead, they are answered in a more mature and scientifically credible way. The work shows where the proposed framework is strongest, where it is mixed, and where future refinements are most needed. In particular, the strongest evidence supports the interpretability and pipeline-design objectives, while the forecasting objective should be interpreted more as enabling structured anomaly analysis than as claiming universal predictive dominance over all baselines.

\subsection{Implications for Future Work}

The discussion also points naturally toward future research directions. The first is to strengthen the forecasting backbone without sacrificing the compactness that makes the dCeNN--ELM design attractive for edge deployment. Since Ridge performs strongly in the current benchmark snapshot, a particularly promising direction would be hybrid or ensemble formulations that preserve the efficiency of linear structure while retaining nonlinear latent features.

The second direction is to deepen the symbolic layer. The current rule base is intentionally compact, which helps interpretability and manageability, but it also limits the semantic breadth of the explanations. Extending the rule base to include richer market logic, uncertainty-aware exceptions, or region-specific operating constraints could make the refinement layer even more valuable.

The third direction is to improve the scalability of the reasoning stage. Incremental ASP, rolling symbolic windows, or multi-tier reasoning architectures could help preserve interpretability while reducing the cost of symbolic refinement in larger or faster environments.

Finally, the finance-aware evaluation could be deepened into more realistic operational or trading scenarios. Even in its current form, it already shifts the evaluation away from abstract anomaly counts and toward downstream usefulness. A stronger future implementation could connect this logic to richer market simulation, probabilistic action policies, or operator decision-support workflows.

\subsection{Closing Synthesis}

In summary, the discussion confirms that the most important contribution of this thesis is not a claim of absolute forecasting superiority, but the demonstration of a coherent and practically motivated neuro-symbolic anomaly-analysis framework. The dCeNN--ELM backbone provides a compact forecasting core; the ASP layer provides interpretability and false-alarm refinement; the software architecture provides modularity and edge-conscious deployment; and the results framework evaluates not only prediction quality, but also anomaly plausibility and downstream usefulness.

That combination is what allows the work to answer its research questions in a meaningful way. The thesis shows that lightweight neuro-symbolic anomaly detection for weather-dependent electricity systems is both feasible and valuable, provided that its strengths and limits are interpreted honestly. In that sense, the research has achieved its central objective: it has demonstrated that combining learning and reasoning can improve the quality of anomaly analysis beyond what either paradigm would provide alone.